\chapter{Introdução}
\label{cap:introducao}

\section{Motivação}

O \mpasa{} (\textit{Model for Prediction Across Scales -- Atmosphere})
\cite{skamarock2012voronoi} é um modelo global/regional de previsão
numérica do tempo cujo núcleo dinâmico opera sobre uma malha horizontal
não estruturada: uma tesselação de Voronoi centroidal esférica (SCVT),
cujos polígonos (predominantemente hexágonos, com um pequeno número de
pentágonos e heptágonos nos pontos de maior curvatura da geração da
malha) definem os volumes de controle das variáveis de massa, enquanto
as arestas entre células vizinhas carregam a componente normal do vento.
Essa escolha de malha -- em contraste com as grades latitude-longitude
regulares tradicionalmente usadas em modelos globais mais antigos --
evita a convergência de meridianos nos polos e permite refinamento local
suave sem as descontinuidades de grades aninhadas.

O projeto no qual este trabalho se insere, \selfgrib{}, nasceu de uma
pergunta simples: é possível gerar a condição inicial e a fronteira
lateral de uma rodada regional do \mpasa{} usando, como fonte de dado
meteorológico, a saída nativa de uma rodada global do \emph{próprio}
\mpasa{}, sem depender de nenhum GRIB externo (GFS, BAM, Eta) e sem
aplicar nenhum \emph{patch} no código-fonte do MPAS-Model? A resposta,
documentada em detalhe no relatório principal do projeto, foi sim: o
\selfgrib{} lê o \texttt{history.nc} de uma rodada global concluída,
extrai e interpola os campos necessários, e escreve o formato binário
intermediário do WPS -- o mesmo formato que o \texttt{ungrib} geraria a
partir de um GRIB real -- permitindo que o \texttt{init\_atmosphere\_model}
original consuma esse dado sem modificação alguma.

Esse desenho tem uma limitação estrutural, no entanto, que motivou
diretamente o trabalho aqui descrito. O caminho ponta a ponta da rota
original é:

\begin{center}
malha nativa global $\to$ \textbf{grade lat-lon regular} $\to$ formato
binário WPS $\to$ malha nativa regional
\end{center}

Duas etapas desse caminho envolvem reamostragem espacial redundante: a
malha nativa global é primeiro remapeada para uma grade lat-lon regular
(pelo utilitário \texttt{convert\_mpas}, da NCAR), e depois o
\texttt{init\_atmosphere\_model} reinterpola essa grade de volta para a
malha nativa \emph{regional} de destino, usando interpolação bilinear ou
de 16 pontos sobre índices de grade estruturada -- o único modo de
interpolação horizontal que esse programa aceita para o caso de dado
real (\texttt{config\_init\_case=7/9}). O resultado é um percurso
\emph{malha $\to$ grade $\to$ malha}: cada etapa de reamostragem introduz
um erro de interpolação de ida e volta, mesmo quando a grade
intermediária é fina, simplesmente pelo fato de trocar de representação
espacial duas vezes em vez de zero. Esse problema já havia se manifestado
concretamente no projeto como uma classe de bug real e documentada (a
grade lat-lon nominal, dimensionada pela elipse teórica do recorte
regional, ficava sistematicamente menor que a extensão real da malha
recortada, deixando células sem dado válido na fronteira -- ver
\cref{cap:discussao}).

\section{Pergunta de pesquisa e escopo do trabalho}

A pergunta que orienta este trabalho é direta: \emph{é possível eliminar
inteiramente a grade lat-lon intermediária, interpolando diretamente
entre a malha de Voronoi de origem e a malha de Voronoi de destino, e
escrevendo os arquivos de entrada do núcleo dinâmico do \mpasa{}
(\texttt{init.nc}/\texttt{lbc.*.nc}) diretamente -- sem passar pelo
formato binário do WPS nem pelo \texttt{init\_atmosphere\_model}
original?}

A investigação começou pela leitura do código-fonte real do
\texttt{init\_atmosphere\_model} (rotina
\texttt{mpas\_init\_atm\_hinterp.F}) para verificar se essa reinterpolação
final para a malha regional aceitaria, de alguma forma, uma lista de
pontos dispersos em vez de uma grade estruturada. A resposta encontrada
foi negativa: o programa é estruturalmente amarrado ao formato de grade
(a interpolação bilinear opera sobre índices \texttt{(i,j)} de uma matriz
regular). Isso significa que \emph{qualquer} ganho de precisão de uma
interpolação nativa malha-a-malha só é alcançável pulando esse programa
inteiramente para a etapa de horizontal/vertical/balanço físico --- não
há um ``escopo mínimo'' que preserve o formato WPS e ainda elimine o
percurso malha$\to$grade$\to$malha. Essa constatação (detalhada na
\cref{cap:arquitetura}) definiu o escopo do trabalho: reimplementar, em
Fortran, cada etapa fisicamente necessária para produzir um
\texttt{init.nc}/\texttt{lbc.*.nc} válido -- não um subconjunto delas.

\section{Contribuições}

Este trabalho contribui com:

\begin{enumerate}[label=(\roman*)]
  \item Um motor de interpolação horizontal nativa malha-a-malha, que
    reaproveita sem modificação o mecanismo de pesos baricêntricos do
    utilitário \texttt{convert\_mpas} (originalmente exposto apenas para
    remapear para grades lat-lon), operando em modo de lista de pontos
    dispersos (\emph{scatter}) sobre os centros de célula da malha de
    destino;
  \item A extração e tradução literal, para um conjunto de módulos
    Fortran independentes e testáveis, das rotinas do núcleo dinâmico do
    MPAS-Model (versão de referência 3.0.2) responsáveis pela geração da
    grade vertical \textit{terrain-following} híbrida, interpolação
    vertical, balanço hidrostático não linear, e inicialização de campos
    de superfície e solo;
  \item Um escritor de \texttt{init.nc}/\texttt{lbc.*.nc} que produz
    arquivos NetCDF binariamente compatíveis com o que o
    \texttt{mpas\_atmosphere} espera, sem depender do
    \texttt{init\_atmosphere\_model} original para gerá-los;
  \item Uma generalização dos programas resultantes para ler os
    parâmetros de configuração (\texttt{config\_*}) do
    \texttt{namelist.init\_atmosphere} real do experimento em tempo de
    execução, tornando a rota reprodutível para qualquer malha e caso de
    estudo, não apenas para o exemplo de validação usado neste relatório;
  \item Validação em duas camadas: (a) numérica, campo a campo, contra
    arquivos \texttt{init.nc}/\texttt{lbc.*.nc} reais de produção; e
    (b) funcional, executando o \texttt{mpas\_atmosphere} real a partir
    dos arquivos gerados por essa rota e obtendo uma previsão de 24 horas
    fisicamente sã, sem erros.
\end{enumerate}

\section{Organização do relatório}

O \cref{cap:fundamentacao} apresenta a fundamentação teórica necessária
-- a estrutura da malha SCVT do \mpasa{}, sua dualidade com a
triangulação de Delaunay, e o princípio da interpolação baricêntrica
usada como método de remapeamento horizontal. O \cref{cap:literatura}
situa esse método na literatura mais ampla sobre malhas de Voronoi em
geociências e computação científica. O \cref{cap:arquitetura} descreve a
arquitetura geral da nova rota (o \emph{pipeline}) e a decisão de
implementação de escopo completo. Os capítulos~\ref{cap:fase1} a
\ref{cap:fase7} descrevem cada fase da implementação em detalhe, com as
equações e algoritmos efetivamente usados. O \cref{cap:implementacao}
discute as escolhas de engenharia da implementação em Fortran. O
\cref{cap:resultados} apresenta os resultados de validação numérica e a
execução funcional do modelo. O \cref{cap:discussao} discute limitações
conhecidas e simplificações deliberadas. O \cref{cap:conclusao} encerra
com as conclusões e sugestões de trabalho futuro.
